% ATLAS-DISTORT2
\section*{ATLAS-DISTORT2: charged coordinate atlases for SAT solution geometry}

\textbf{STATUS: COMPUTATIONAL finite experiment plus KNOWN-THEOREM model correspondence. P versus NP remains OPEN.}

\paragraph{Repair being tested.}
SOLUTION-GEOMETRY1 falsified raw shattering, freezing, sparsity, and single-flip distortion as universal hardness measures. The proposed repair is to measure geometry only after removing structure captured by an explicitly declared and efficiently acquired tractable atlas.

The acquisition qualifier is load-bearing. A singleton solution set is mathematically affine, but an arbitrary CNF is not granted a Gaussian chart merely because exhaustive search later revealed one solution. The chart must be recovered from the same input by the declared recognition algorithm.

\paragraph{Atlas 1: affine/XOR coordinates.}
Standard 3-XOR constraints are encoded by the four 3-clauses excluding assignments of the wrong parity. The experiment syntactically recognizes complete four-clause XOR blocks and then applies Gaussian elimination. If the parity matrix has rank $n-d$, the solution set is an affine translate of $\mathbb F_2^d$.

A hostile affine control was constructed by choosing a random four-dimensional nullspace and many weight-three parity equations spanning its orthogonal complement. The resulting CNFs have exactly 16 satisfying assignments but exact primal pathwidth 10--12 at $n=20$.

Before normalization these instances are maximally misleading geometrically: all 16 satisfying assignments are isolated under one-bit flips, every original variable is frozen within its singleton component, and the disconnected-pair fraction is one. Nevertheless the XOR blocks are directly recognized, Gaussian elimination exposes four free coordinates, and the normalized solution space is exactly
\[ \mathbb F_2^4. \]
Thus in Gaussian coordinates it has one component, zero frozen directions, unit local flexibility, unit intrinsic/Hamming stretch, and no residual affine constraints.

\paragraph{Atlas 2: bounded pathwidth.}
For the local-chain control, clauses involve only consecutive triples. Exact pathwidth was computed through the equivalent vertex-separation formulation. Kinnersley proves that vertex separation equals pathwidth.

At $n=20$ the two chain controls have raw disconnected-pair fractions approximately $0.645$ and $0.755$, with 3 and 4 raw components and 10 and 14 frozen variables in the largest component. Yet their exact primal pathwidth is
\[ \boxed{2}. \]
Along an optimal pathwidth ordering, only 4 and 3 extendible separator assignments are ever simultaneously needed. Thus the natural dynamic-programming atlas replaces severe raw fragmentation by a layered representation with at most two bits of interface state.

\paragraph{Random 3-SAT comparison.}
For satisfiable random simple genuine 3-SAT at $n=20$, the exact primal pathwidths in this finite panel are
\[
\begin{array}{c|c}
\alpha & \text{pathwidths}\\
\hline
3.0&12,12\\
3.5&14,13\\
4.0&15,15\\
4.25&14,13.
\end{array}
\]
The corresponding maxima of extendible separator projections are
\[ 139,81;\quad 47,51;\quad 14,22;\quad 14,1. \]
No tested random formula is recognized as a complete XOR-block system.

The final value $1$ is an important negative control rather than evidence of easiness. That instance has a single satisfying assignment, so after exhaustive semantic knowledge every separator has only one extendible state. But its primal pathwidth is 13 and no Gaussian chart is syntactically available. This demonstrates that semantic state count after solving cannot be treated as free acquisition.

\paragraph{What the experiment supports.}
The right object is therefore relative to a fixed atlas class $\mathcal A$:
\[
\operatorname{Dist}_{\mathcal A}(F)
=
\min_{R\in\mathcal A(F)}
\operatorname{ResidualObstruction}(F/R),
\]
where membership of $R$ in $\mathcal A(F)$ means the representation is actually recovered from $F$ by the declared algorithm and its acquisition cost is charged.

In the present atlas,
\[
\mathcal A=
\{\text{standard XOR blocks + Gaussian elimination},\
\ \text{bounded-pathwidth decomposition + dynamic programming}\}.
\]
The affine and low-pathwidth hostile controls collapse as intended, while the tested random 3-SAT instances do not receive either chart.

\paragraph{Limitations.}
This remains a CANDIDATE architecture, not a lower bound. An unknown polynomial-time representation may flatten formulas outside the current atlas. Worse, minimization over all efficient coordinate systems is not legitimate unless the class and acquisition procedure are fixed independently; otherwise the coordinate search can hide the SAT solver itself. Finally, no theorem currently transfers large residual normalized geometry to arbitrary SAT time.

\paragraph{Next target.}
Add Horn and 2-SAT charts, then replace component counts by continuation incompatibility across the acquired chart interfaces. The strongest candidate would be a single residual state/cut quantity that is small after Gaussian, Horn, 2-SAT, and bounded-width normalizations but remains large on the random hard region and admits a theorem-level connection to branching-program or proof-complexity resources.

\begin{thebibliography}{9}
\bibitem{Kinnersley92} N.~G.~Kinnersley, \emph{The Vertex Separation Number of a Graph Equals Its Path-Width}, Information Processing Letters 42(6):345--350, 1992. doi:10.1016/0020-0190(92)90234-M.
\bibitem{GKMP09} P.~Gopalan, P.~G.~Kolaitis, E.~Maneva, and C.~H.~Papadimitriou, \emph{The Connectivity of Boolean Satisfiability: Computational and Structural Dichotomies}, SIAM J. Comput. 38(6):2330--2355, 2009. doi:10.1137/07070440X.
\bibitem{SamerSzeider10} M.~Samer and S.~Szeider, \emph{A Fixed-Parameter Algorithm for \#SAT with Parameter Incidence Treewidth}, J. Comb. Optim. 19:100--121, 2010.
\bibitem{Lampis24} M.~Lampis, \emph{The Primal Pathwidth SETH}, arXiv:2403.07239, 2024.
\end{thebibliography}

% END_OF_NOTE::ATLAS-DISTORT2